\section{Resultados e Discussão}

\begin{frame}{Resultados e Discussão}
  Os painéis (solução vs.\ referência, erro relativo e espectro de Fourier) e as animações
  de cada método são apresentados \alert{ao vivo}. A seguir, os principais comentários por
  arquitetura.

  \vspace{0.3cm}
  \begin{information}[O que observar]
    Em todos os casos, o comportamento de interesse aparece ao \textbf{aumentar} $\alpha=\beta$:
    é quando a solução espalha energia por mais frequências e o \alert{viés espectral} e o
    \alert{efeito de Gibbs} se tornam visíveis.
  \end{information}
\end{frame}

\begin{frame}{FNO}
  \begin{itemize}
    \item Inferência \alert{muito rápida} --- uma única passada direta, mesmo com
          $\approx 1$M de parâmetros;
    \item boa fidelidade nos regimes vistos no treino;
    \item porém surge o \alert{efeito de Gibbs} na direção temporal: oscilações artificiais
          que não desaparecem aumentando o número de modos.
  \end{itemize}
\end{frame}

\begin{frame}{PINO --- com dados vs.\ sem dados}
  \begin{itemize}
    \item \textbf{Com dados:} o resíduo físico \alert{elimina o Gibbs} e a aproximação fica
          consistente com a referência;
    \item a \textbf{propagação da condição inicial} no tempo é o ponto mais delicado do
          treino;
    \item \textbf{Sem dados:} reaparece um \alert{viés espectral forte} --- o PINO puramente
          físico tem dificuldade com os detalhes de alta frequência.
  \end{itemize}
\end{frame}

\begin{frame}{PINN e o viés espectral}
  \begin{itemize}
    \item A PINN \alert{sofre viés espectral}: aprende primeiro as componentes suaves e
          perde as altas frequências do perfil;
    \item acrescentar um termo de \textbf{dados} \emph{atenua} (mas não elimina) o efeito;
    \item \textbf{sem dados}, a dificuldade é mais severa à medida que a não linearidade cresce.
  \end{itemize}

  \begin{information}[Comparando as arquiteturas]
    Por aprenderem diretamente no espaço de Fourier, o \textbf{FNO} e o \textbf{PINO com
    dados} sofrem menos com o viés espectral; ele reaparece com força no \textbf{PINO sem
    dados} e nas \textbf{PINNs}.
  \end{information}
\end{frame}
